% Part: proof-theory
% Chapter: sequent-calculus
% Section: rules-G1c

\documentclass[../../../include/open-logic-section]{subfiles}

\begin{document}

\begin{table}
\begin{tabular}{@{}c@{}c@{}}
\multicolumn{2}{@{}c@{}}{అక్షయాలు}\\[2ex]
$!A \Sequent !A$ & $\lfalse \Sequent \quad$
\\[2ex]
\multicolumn{2}{@{}c@{}}{నిర్మాణ నియమాలు}\\[2ex]
\Axiom$ \Gamma \fCenter \Delta $
\RightLabel{\LeftR{\Weakening}}
\UnaryInf$ !A, \Gamma \fCenter \Delta$
\DisplayProof
&
\Axiom$ \Gamma \fCenter $
\RightLabel{\RightR{\Weakening}}
\UnaryInf$ \Gamma \fCenter !A$
\DisplayProof
\\[4ex]
\Axiom$ !A, !A, \Gamma \fCenter \Delta $
\RightLabel{\LeftR{\Contraction}}
\UnaryInf$ !A, \Gamma \fCenter \Delta$
\DisplayProof
&
\\[4ex]
\multicolumn{2}{@{}c@{}}{తార్కిక నియమాలు}\\[2ex]
\Axiom$ \Gamma \fCenter !A $
\RightLabel{\LeftR{\lnot}}
\UnaryInf$ \lnot !A, \Gamma \fCenter $
\DisplayProof
&
\Axiom$!A, \Gamma \fCenter $
\RightLabel{\RightR{\lnot}}
\UnaryInf$ \Gamma \fCenter \lnot !A $
\DisplayProof
\\[4ex]\begin{tabular}[t]{@{}l@{}}
\Axiom$ !A, \Gamma \fCenter \Delta$
\RightLabel{\LeftR{\land}}
\UnaryInf$ !A \land !B, \Gamma \fCenter \Delta$
\DisplayProof
\\[3ex]
\Axiom$!B, \Gamma \fCenter \Delta$
\RightLabel{\LeftR{\land}}
\UnaryInf$!A \land !B, \Gamma \fCenter \Delta$
\DisplayProof\\[3ex]
\end{tabular}
&
\Axiom$\Gamma \fCenter !A$
\Axiom$ \Gamma \fCenter !B$
\RightLabel{\RightR{\land}}
\BinaryInf$ \Gamma \fCenter !A \land !B $
\DisplayProof
\\[4ex]\Axiom$!A, \Gamma \fCenter \Delta$
\Axiom$!B, \Gamma \fCenter \Delta$
\RightLabel{\LeftR{\lor}}
\BinaryInf$!A \lor !B, \Gamma \fCenter \Delta$
\DisplayProof
&
\begin{tabular}[t]{@{}r@{}}
\Axiom$\Gamma \fCenter !A$
\RightLabel{\RightR{\lor}}
\UnaryInf$ \Gamma \fCenter !A \lor !B$
\DisplayProof
\\[3ex]
\Axiom$ \Gamma \fCenter !B$
\RightLabel{\RightR{\lor}}
\UnaryInf$ \Gamma \fCenter !A \lor !B$
\DisplayProof\\[3ex]
\end{tabular}
\\[4ex]
\Axiom$ \Gamma \fCenter !A$
\Axiom$ !B, \Gamma \fCenter \Delta$
\RightLabel{\LeftR{\lif}}
\BinaryInf$ !A \lif !B, \Gamma \fCenter \Delta$
\DisplayProof
&
\Axiom$ !A, \Gamma \fCenter !B$
\RightLabel{\RightR{\lif}}
\UnaryInf$ \Gamma \fCenter !A \lif !B $
\DisplayProof
\\[4ex]
\Axiom$ !A(t), \Gamma \fCenter \Delta$
\RightLabel{\LeftR{\lforall}}
\UnaryInf$ \lforall[x][!A(x)],\Gamma \fCenter \Delta$
\DisplayProof
&
\Axiom$ \Gamma \fCenter !A(a) $
\RightLabel{\RightR{\lforall}}
\UnaryInf$ \Gamma \fCenter \lforall[x][!A(x)]$
\DisplayProof
\\[4ex]
\Axiom$ !A(a), \Gamma \fCenter \Delta $
\RightLabel{\LeftR{\lexists}}
\UnaryInf$ \lexists[x][!A(x)], \Gamma \fCenter \Delta$
\DisplayProof
&
\Axiom$ \Gamma \fCenter !A(t) $
\RightLabel{\RightR{\lexists}}
\UnaryInf$ \Gamma \fCenter \lexists[x][!A(x)]$
\DisplayProof
\end{tabular}
\caption{అంతఃప్రజ్ఞావాద తర్కానికి \Log{G1i} నియమాలు. $\RightR{\forall}$, $\LeftR{\exists}$లలో \tetoken{స్థిరాంకం}{!!{constant}}~$a$ నిర్ధారణ సీక్వెంట్‌లో రాకూడదు. $\Gamma$ \tetoken{సూత్రాల}{!!{formula}s} బహుసమితి; $\Delta$ ఖాళీగా లేదా ఒకే \tetoken{సూత్రంతో}{!!{formula}} ఉండవచ్చు. \RightR{\Contraction} నియమం లేదు. \RightR{\Weakening}ను, అక్షయం $\lfalse \Sequent \quad$ను తీసేసిన \Log{G1i} వ్యవస్థ \Log{G1m}.}
\ollabel{tab:G1i}
\end{table}
\sourcecorrection{OLTEPTSEQG1I-001}{మూల అంతఃప్రజ్ఞావాద పట్టికకు సాంప్రదాయ పట్టిక గుర్తే ఇచ్చారు; ఈ వ్యవస్థకు ప్రత్యేక పట్టిక గుర్తును ఇచ్చాం.}
\sourcecorrection{OLTEPTSEQG1I-002}{కనిష్ఠ వ్యవస్థ వివరణలో మూలం పక్క సందర్భాలున్న అసత్య అక్షయ కుటుంబాన్ని పేర్కొంది; ఈ పట్టికలోని అసలు ఆధార అసత్య అక్షయాన్ని నిర్దిష్టంగా రాశాం.}

\end{document}
